\documentclass[11pt,a4paper]{article}
\usepackage[margin=1in]{geometry}
\usepackage{amsmath,booktabs,hyperref,float}

\title{Iteration 028: Huber + Correlation Loss}
\author{Autoresearch Loop}
\date{April 8, 2026}

\begin{document}
\maketitle

\begin{abstract}
Replacing MSE with Huber loss in the combined objective, with InstanceNorm,
achieves $r=0.375$. Huber's robustness to outliers does not compensate for
InstanceNorm's information loss.
\end{abstract}

\section{Hypothesis}
The MSE component is sensitive to outliers from noisy channels (ELC, ERT).
Huber loss (L2 for small errors, L1 for large) reduces outlier influence
while maintaining gradient stability near zero error.

\section{Method}
Loss: $0.5 \times \text{Huber}(\delta=1.0) + 0.5 \times (1-r)$. Architecture:
InstanceNorm + FIR (7 taps). Corr validation + 15\% channel dropout.

\section{Results}
\begin{table}[H]
\centering
\begin{tabular}{lrrr}
\toprule
Model & Mean $r$ & Std $r$ & SNR (dB) \\
\midrule
Combined loss (iter017) & 0.378 & 0.078 & 0.61 \\
Huber+corr + InstanceNorm & 0.375 & 0.075 & 0.62 \\
\bottomrule
\end{tabular}
\end{table}

\section{Analysis}
The lowest std of any model tested (0.075 vs typical 0.077) suggests Huber
loss does improve consistency. However, InstanceNorm's $-$0.003 effect
dominates. Huber loss without InstanceNorm is worth testing.

\section{Next}
Focus on non-InstanceNorm variants to isolate technique effects.

\end{document}
